% A2 — Abstract.  DRAFT 1, 6 Oct 2026 (~150 words).
Liu et al.\ explain grokking with an effective theory that predicts a critical
training data fraction $r_c \approx 0.4$ for a toy addition task. I test this prediction with the authors' released code. The training-free part reproduces. The probability that the training data determine the
linear representation crosses one half at $r = 0.41$, and my implementation agrees with the authors' on all 1800 training sets checked. Trained networks,
however, do not follow the prediction with the released settings. Over twenty
seeds per training size, half of the runs succeed within $10^4$ steps only at
$r \approx 0.76$. Only 70 of the 138 runs whose data determine the representation actually find it.
Decoder weight decay has almost no effect. The decoder learning rate matters
a lot, since a slower decoder raises this number to 96 and a faster one
lowers it to 15. Whether a network finds the structure thus depends on how
fast the decoder learns, which the theory does not model.
